\scanpage{135}
\mspage{6}
% archive sheet number: 134

\[
\text{c)}\quad X'\mathrel{|\circ|}\varnothing\qquad\text{tjs satisfaite.}
\]

L'introduction de $\widehat{A}'$, et les conventions précédentes, sont inspirées par l'exemple $2^\circ)$ précédent.

\uline{Noyau} de $f$ :
\[
\operatorname{Ker} f\overset{\mathrm{dfn}}{=}\{X\in A\mid f(X)=\varnothing\}
\]
Alors b) implique que $\operatorname{Ker} f\subset A$ est fermé pour $\overset{\circ}{\ll}$. Du coup on s'attend que ce soit dans $\Sigma_A$, mais on doit sans doute pour cela imposer des conditions \unsure{supplémentaires} sur $f$ (en tous cas c'est OK dans l'exemple $2^\circ$ et ses variantes).

Partons d'une partie $\unsure{N}\subset A$, fermée pour $\overset{\circ}{\ll}$, est-ce le noyau d'un hom. d'atelier, et ce dernier peut-il se choisir canoniquement ? Plus précisément, peut-on trouver un atelier $A/N$ (opération d'\uline{excision})
